\uuid{FmEz}
\exo7id{3555}
\titre{exo7 3555}
\auteur{quercia}
\organisation{exo7}
\datecreate{2010-03-10}
\isIndication{false}
\isCorrection{false}
\chapitre{Réduction d'endomorphisme, polynôme annulateur}
\sousChapitre{Polynôme annulateur}
\cpge{mp-mpi.reduction-endomorphismes.d, pc.reduction.c, psi.reduction-endomorphismes.c, pt.reduction-endomorphismes.c, mp-mpi.reduction-endomorphismes.g, pc.complements-espaces-vectoriels.d, psi.complements-espaces-vectoriels.d, mp-mpi.reduction-endomorphismes.i, pc.reduction.d, psi.reduction-endomorphismes.d}
\module{Algèbre}
\niveau{L2}
\difficulte{}

\contenu{
\texte{
Soit $E$ un $ K$-ev de dimension finie et $u \in \mathcal{L}(E)$.
}
\begin{enumerate}
    \item \question{On suppose $u$  diagonalisable et on note
    $\lambda_1,\dots,\lambda_p$ ses valeurs propres distinctes.
  }
    \item \question{Montrer qu'il existe des endomorphismes $u_1,\dots,u_p$ tels que pour tout
    polynôme $P \in  K[X]$, on ait :\par
    $P(u) = \sum_{i=1}^p P(\lambda_i) u_i$.}
    \item \question{Montrer qu'il existe un polynôme $P_i$ tel que $u_i = P_i(u)$.}
\end{enumerate}
}
